\documentclass[11pt,a4paper]{article}
\usepackage[utf8]{inputenc}
\usepackage[indonesian]{babel}
\usepackage{amsmath,amsfonts,amssymb,amsthm}
\usepackage{graphicx}
\usepackage{geometry}
\usepackage{fancyhdr}
\usepackage{tcolorbox}
\usepackage{booktabs}
\usepackage{tabularx}
\usepackage{tikz}
\usetikzlibrary{positioning,shapes.geometric,arrows.meta,calc}
\usepackage{circuitikz}
\usepackage{enumitem}
\usepackage{listings}
\usepackage{xcolor}
\usepackage{hyperref}

\geometry{top=2.5cm, bottom=2.5cm, left=2.5cm, right=2.5cm, headheight=15pt}

\pagestyle{fancy}
\fancyhf{}
\fancyhead[L]{\small\textbf{Persamaan Diferensial 2026} -- Teknik Elektro UNIB}
\fancyhead[R]{\small Modul 7: Invers Laplace \& Analisis Rangkaian Domain $s$}
\fancyfoot[C]{\thepage}

% Pengaturan kode Julia
\lstdefinelanguage{Julia}%
  {morekeywords={abstract,break,case,catch,const,continue,do,else,elseif,%
      end,export,false,for,function,global,if,import,%
      let,local,macro,module,mutable,primitive,quote,return,struct,%
      true,try,type,using,var,while},%
   sensitive=true,%
   morecomment=[l]{\#},%
   morecomment=[n]{\#=}{=\#},%
   morestring=[b]',%
   morestring=[b]"%
  }[keywords,comments,strings]

\lstdefinestyle{juliastyle}{
    language=Julia,
    basicstyle=\footnotesize\ttfamily,
    keywordstyle=\color{blue!80!black}\bfseries,
    commentstyle=\color{green!50!black}\itshape,
    stringstyle=\color{red!80!black},
    numbers=left,
    numberstyle=\tiny\color{gray},
    breaklines=true,
    showstringspaces=false,
    frame=single,
    backgroundcolor=\color{gray!6},
    captionpos=b,
    xleftmargin=10pt,
    tabsize=4
}

\definecolor{headercolor}{RGB}{0,51,102}
\definecolor{accentcolor}{RGB}{0,102,204}
\definecolor{shadecolor}{RGB}{245,248,253}

\hypersetup{
    colorlinks=true,
    linkcolor=headercolor,
    citecolor=accentcolor,
    urlcolor=blue
}

\theoremstyle{definition}
\newtheorem{definition}{Definisi}[section]
\newtheorem{example}{Contoh Soal}[section]
\newtheorem{theorem}{Teorema}[section]

\title{\textbf{\Large MODUL AJAR KULIAH MINGGU 7}\\[0.3cm]
\textbf{\LARGE INVERS TRANSFORMASI LAPLACE, EKSPANSI PECAHAN PARSIAL, ANALISIS RANGKAIAN DOMAIN $s$, DAN TEOREMA NILAI AWAL/AKHIR}}
\author{\textbf{Ir. Novalio Daratha, S.T., M.Sc., Ph.D.} \& \textbf{Muhammad Arfan, S.T., M.T.} \\ 
Jurusan Teknik Elektro -- Fakultas Teknik \\ 
Universitas Bengkulu}
\date{Semester Genap 2026}

\begin{document}

\maketitle
\begin{center}
  \vspace{-0.25cm}
  \begin{tcolorbox}[colback=red!4!white,colframe=red!75!black,boxrule=0.5pt,arc=2pt,left=6pt,right=6pt,top=3pt,bottom=3pt,width=0.98\textwidth]
    \centering\footnotesize
    \textbf{Video Perkuliahan Resmi (YouTube):} {\color{red!80!black}\textbf{\url{https://youtu.be/crWRyniSoi0}}} \quad $\bullet$ \quad \textbf{Playlist:} \href{https://www.youtube.com/playlist?list=PLS5oOZWXZeTw}{\color{blue!80!black}\textbf{bit.ly/pd-unib-playlist}} \quad $\bullet$ \quad \textbf{Portal:} \url{https://pd.ndaratha.my.id}
  \end{tcolorbox}
  \vspace{-0.15cm}
\end{center}


\begin{tcolorbox}[colback=blue!5!white,colframe=headercolor,title=\textbf{Capaian Pembelajaran Khusus (Sub-CPMK 4)}]
Setelah menyelesaikan pembelajaran pada Modul Ajar Minggu 7 ini, mahasiswa diharapkan mampu:
\begin{enumerate}[leftmargin=*]
    \item \textbf{C1 (Mengingat):} Menyatakan definisi formal Invers Transformasi Laplace melalui integral kontur Bromwich, sifat kelinieran inversi, serta bentuk-bentuk kanonik pecahan parsial untuk kutub real sederhana, berulang, dan kompleks konjugat.
    \item \textbf{C2 (Memahami):} Menjelaskan konsep impedansi dan admitansi operasional di domain frekuensi kompleks $s$, serta makna fisis pembentukan model sumber semu (\textit{fictitious initial-condition sources}) $L i(0^-)$ dan $v_C(0^-)/s$ yang merepresentasikan memori energi tersimpan pada elemen induktor dan kapasitor.
    \item \textbf{C3 (Menerapkan):} Menghitung dekomposisi pecahan parsial fungsi aljabar rasional $F(s)$ menggunakan metode penutupan Heaviside (\textit{cover-up method}) dan derivasi bertingkat untuk mengembalikan fungsi ke domain waktu fisis $f(t)$.
    \item \textbf{C4 (Menganalisis):} Menganalisis respon transien rangkaian dinamis linier orde tinggi secara simultan di domain $s$ menggunakan Hukum Kirchoff (KCL/KVL), serta memisahkan solusi total menjadi Respon Masukan-Nol (\textit{Zero-Input Response} / ZIR) dan Respon Keadaan-Nol (\textit{Zero-State Response} / ZSR).
    \item \textbf{C5 (Mengevaluasi):} Mengevaluasi perilaku transien sesaat ($t \to 0^+$) dan keadaan tunak (\textit{steady-state}, $t \to \infty$) tanpa memerlukan inversi penuh menggunakan Teorema Nilai Awal (\textit{Initial Value Theorem} / IVT) dan Teorema Nilai Akhir (\textit{Final Value Theorem} / FVT) pada fenomena tegangan pemulihan transien (\textit{Transient Recovery Voltage} / TRV) pemutus tenaga tegangan tinggi (IEEE C37.011 / IEC 62271-100).
    \item \textbf{C6 (Menciptakan/Komputasi):} Mengembangkan program komputasi ilmiah terbuka berbasis \textbf{Julia} (\texttt{DifferentialEquations.jl} \& \texttt{Plots.jl}) untuk mensimulasikan dinamika transien sistem RLC orde dua berenergi awal dan memvalidasi kecocokan solusi analitis invers Laplace terhadap metode numerik adaptif.
\end{enumerate}
\end{tcolorbox}

\vspace{0.3cm}
\tableofcontents
\vspace{0.5cm}

\newpage

\section{Peta Konsep \& Posisi Modul dalam Kurikulum}

Modul Minggu 7 merupakan puncak kulminasi dari Paruh Pertama kurikulum Persamaan Diferensial (PDB dan Rangkaian Listrik). Setelah pada Minggu 6 kita mempelajari bagaimana mentransformasikan sinyal waktu dan operator turunan ke dalam ranah aljabar frekuensi kompleks $s$, kini kita menutup siklus operasional tersebut dengan mempelajari teknik rekonstruksi sinyal kembali ke ranah waktu fisis: \textbf{Invers Transformasi Laplace} ($\mathcal{L}^{-1}$).

Melalui penguasaan Invers Laplace dan konsep pemodelan rangkaian ekivalen domain $s$, mahasiswa Teknik Elektro tidak perlu lagi menurunkan persamaan diferensial orde tinggi di domain waktu menggunakan kalkulus yang melelahkan. Seluruh hukum dasar rangkaian (Hukum Ohm, KVL, KCL, Teorema Thevenin/Norton) dapat langsung diaplikasikan di domain $s$ layaknya menganalisis rangkaian arus searah (DC), di mana syarat awal energi terakomodasi secara otomatis sebagai sumber tegangan atau arus semu.

Posisi strategis Modul Minggu 7 disajikan pada diagram kurikulum Gambar~\ref{fig:peta_jalan_pd_w7}.

\begin{figure}[htbp]
\centering
\resizebox{\linewidth}{!}{%
\begin{tikzpicture}[
    node distance=0.85cm and 0.45cm,
    block/.style={rectangle, draw=headercolor, fill=blue!8, thick, rounded corners=3pt, align=center, text width=3.35cm, minimum height=1.05cm, font=\scriptsize},
    doneblock/.style={rectangle, draw=green!60!black, fill=green!10, thick, rounded corners=3pt, align=center, text width=3.35cm, minimum height=1.05cm, font=\scriptsize},
    highlight/.style={rectangle, draw=red!80!black, fill=red!15, very thick, rounded corners=3pt, align=center, text width=3.35cm, minimum height=1.05cm, font=\scriptsize\bfseries},
    midblock/.style={rectangle, draw=orange!80!black, fill=orange!15, thick, rounded corners=3pt, align=center, text width=3.35cm, minimum height=1.05cm, font=\scriptsize\bfseries},
    arrow/.style={->, >=stealth, line width=1.1pt, headercolor}
]
    % Paruh 1: PDB & Rangkaian Listrik
    \node[doneblock] (w1) {Minggu 1: Klasifikasi PD, IVP\\\& Pemodelan Fisis RL};
    \node[doneblock, right=of w1] (w2) {Minggu 2: PDB Orde 1\\(Separabel \& Faktor Integrasi)};
    \node[doneblock, right=of w2] (w3) {Minggu 3: Aplikasi Rekayasa\\(Transien RC, RL, Termal)};
    \node[doneblock, right=of w3] (w4) {Minggu 4: PDB Orde 2 Homogen\\(Wronskian \& Tangki LC)};
    
    \node[doneblock, below=of w4] (w5) {Minggu 5: PDB Orde 2 Non-Homogen\\\& Transien RLC Seri AC};
    \node[doneblock, left=of w5] (w6) {Minggu 6: Transformasi Laplace\\Dasar \& Teorema Geser};
    \node[highlight, left=of w6] (w7) {Minggu 7: Invers Laplace \&\\Analisis Rangkaian Domain $s$};
    \node[midblock, left=of w7] (w8) {Minggu 8: Evaluasi Capaian\\Ujian Tengah Semester (UTS)};
    
    % Paruh 2: PDP & Medan Elektromagnetika
    \node[block, below=of w8] (w9) {Minggu 9: Pengantar PDP \&\\Analisis Deret Fourier};
    \node[block, right=of w9] (w10) {Minggu 10: Solusi PDP Metode\\Pemisahan Variabel};
    \node[block, right=of w10] (w11) {Minggu 11: Persamaan Gelombang\\1D (Telegrafer Transmisi)};
    \node[block, right=of w11] (w12) {Minggu 12: Persamaan Panas 1D\\\& Manajemen Termal Kabel};
    
    \node[block, below=of w12] (w13) {Minggu 13: Persamaan Laplace 2D\\Potensial Elektrostatik};
    \node[block, left=of w13] (w14) {Minggu 14: Fungsi Bessel \&\\Efek Kulit (\textit{Skin Effect})};
    \node[block, left=of w14] (w15) {Minggu 15: 4 Persamaan Maxwell\\Diferensial \& Sintesis PDP};
    \node[midblock, left=of w15] (w16) {Minggu 16: Evaluasi Akhir\\Ujian Akhir Semester (UAS)};
    
    \draw[arrow] (w1) -- (w2);
    \draw[arrow] (w2) -- (w3);
    \draw[arrow] (w3) -- (w4);
    \draw[arrow] (w4) -- (w5);
    \draw[arrow] (w5) -- (w6);
    \draw[arrow] (w6) -- (w7);
    \draw[arrow] (w7) -- (w8);
    \draw[arrow] (w8) -- (w9);
    \draw[arrow] (w9) -- (w10);
    \draw[arrow] (w10) -- (w11);
    \draw[arrow] (w11) -- (w12);
    \draw[arrow] (w12) -- (w13);
    \draw[arrow] (w13) -- (w14);
    \draw[arrow] (w14) -- (w15);
    \draw[arrow] (w15) -- (w16);
\end{tikzpicture}%
}
\caption{Peta jalan kurikulum perkuliahan dengan penekanan pada Modul Minggu 7 sebagai penutup Paruh 1.}
\label{fig:peta_jalan_pd_w7}
\end{figure}

\section{Pendahuluan \& Filosofi Rekayasa: Dari Ranah Aljabar Kembali ke Ranah Waktu}

Pada proses pemecahan masalah keteknikan menggunakan kalkulus operasional, kita menempuh rute tiga langkah yang sangat sistematis:
\begin{enumerate}[leftmargin=*]
    \item \textbf{Pemetaan Maju (\textit{Forward Transform}):} Masalah fisik nyata yang dinyatakan dalam bentuk persamaan diferensial integro-diferensial domain waktu dipetakan ke domain aljabar $s$ melalui Transformasi Laplace.
    \item \textbf{Penyelesaian Aljabar (\textit{Algebraic Manipulation}):} Di domain $s$, operasi turunan berubah menjadi perkalian $s$ dan operasi integral berubah menjadi pembagian $s$. Kita menyelesaikan sistem persamaan aljabar linier secara simultan untuk memperoleh variabel respon target, misalnya $V(s)$ atau $I(s)$.
    \item \textbf{Rekonstruksi Balik (\textit{Inverse Transform}):} Fungsi aljabar rasional $F(s)$ yang diperoleh harus dikembalikan ke dunia nyata menjadi fungsi waktu $f(t)$ agar rekayasawan dapat membaca bentuk gelombang tegangan, lonjakan arus transien, atau kestabilan sistem.
\end{enumerate}

Tantangan utama pada tahap ketiga adalah bahwa hasil manipulasi aljabar sering kali berupa rasio polinomial berderajat tinggi:
\begin{equation}
    F(s) = \frac{N(s)}{D(s)} = \frac{b_m s^m + b_{m-1}s^{m-1} + \dots + b_1 s + b_0}{a_n s^n + a_{n-1}s^{n-1} + \dots + a_1 s + a_0}
    \label{eq:rational_transfer}
\end{equation}
Bentuk rasional semacam ini umumnya tidak tertera langsung dalam tabel baku pasangan transformasi Laplace. Oleh karena itu, kita memerlukan jembatan aljabar yang memecah rasio kompleks tersebut menjadi penjumlahan suku-suku pecahan sederhana yang masing-masing berkorespondensi langsung dengan fungsi waktu elementer (eksponensial, sinus, kosinus, atau ramp). Teknik matematis inilah yang disebut \textbf{Ekspansi Pecahan Parsial} (\textit{Partial Fraction Expansion} / PFE).

\section{Definisi Formal Invers Laplace \& Teorema Integral Bromwich}

\subsection{Definisi Formal Integral Kontur Bromwich}

\begin{definition}[Invers Transformasi Laplace]
Secara analitis murni, jika $F(s) = \mathcal{L}\{f(t)\}$ memiliki wilayah konvergensi (\textit{Region of Convergence}) pada setengah bidang $\text{Re}\{s\} > \gamma_0$, maka invers transformasi Laplace dari $F(s)$ dinotasikan dengan $\mathcal{L}^{-1}\{F(s)\} = f(t)$, dan ditentukan oleh \textbf{Integral Kontur Kompleks Bromwich} (\textit{Bromwich Inversion Integral}):
\begin{equation}
    f(t) = \mathcal{L}^{-1}\{F(s)\} = \frac{1}{2\pi j} \int_{\gamma - j\infty}^{\gamma + j\infty} F(s) e^{st} \, ds, \quad \text{untuk } t > 0
    \label{eq:bromwich_integral}
\end{equation}
di mana garis integrasi vertikal $s = \gamma$ terletak di sebelah kanan seluruh titik singularitas (\textit{poles}) dari fungsi $F(s)$ di bidang kompleks $s$ ($\gamma > \gamma_0$).
\end{definition}

Berdasarkan Teorema Residu Cauchy dari analisis variabel kompleks, integral Bromwich~\eqref{eq:bromwich_integral} setara dengan penjumlahan residu pada semua kutub (\textit{poles}) dari fungsi $F(s)e^{st}$:
\begin{equation}
    f(t) = \sum \text{Residue} \left\{ F(s) e^{st} \right\}_{\text{pada seluruh kutub } p_k}
    \label{eq:cauchy_residue}
\end{equation}

Meskipun integral kontur dan kalkulus residu memberikan landasan matematis yang sangat kuat, dalam praktik rekayasa elektro kita menggunakan pendekatan operasional yang jauh lebih praktis dan elegan: \textbf{Tabel Transformasi Baku yang dikombinasikan dengan Ekspansi Pecahan Parsial}.

\subsection{Sifat Kelinieran Invers Transformasi Laplace}

\begin{theorem}[Kelinieran Invers Laplace]
Sama seperti transformasi maju, operator invers transformasi Laplace $\mathcal{L}^{-1}$ adalah operator linier. Jika $F_1(s)$ dan $F_2(s)$ memiliki invers berturut-turut $f_1(t)$ dan $f_2(t)$, serta $c_1, c_2$ adalah konstanta sebarang, maka:
\begin{equation}
    \mathcal{L}^{-1}\left\{ c_1 F_1(s) + c_2 F_2(s) \right\} = c_1 \mathcal{L}^{-1}\{F_1(s)\} + c_2 \mathcal{L}^{-1}\{F_2(s)\} = c_1 f_1(t) + c_2 f_2(t)
    \label{eq:inverse_linearity}
\end{equation}
\end{theorem}

Sifat inilah yang menjamin bahwa jika kita dapat memecah fungsi rasional rumit $F(s)$ menjadi penjumlahan pecahan-pecahan sederhana, kita cukup mencari invers dari tiap pecahan sederhana tersebut secara terpisah lalu menjumlahkannya.

\section{Metode Ekspansi Pecahan Parsial (Partial Fraction Expansion)}

Diberikan fungsi rasional sejati (\textit{strictly proper rational function}) $F(s) = \frac{N(s)}{D(s)}$ di mana derajat pembilang lebih kecil dari derajat penyebut ($\text{deg}(N) < \text{deg}(D) = n$). Jika pembilang memiliki derajat lebih besar atau sama ($\text{deg}(N) \ge \text{deg}(D)$), pembagian polinomial panjang harus dilakukan terlebih dahulu hingga menghasilkan suku polinomial murni ditambah fungsi rasional sejati.

Akar-akar dari penyebut $D(s) = 0$ disebut sebagai \textbf{kutub-kutub} (\textit{poles}) dari $F(s)$. Terdapat tiga kasus mendasar bentuk kutub yang umum dijumpai:

\subsection{Kasus 1: Kutub-Kutub Real Sederhana (Distinct Real Poles)}
Jika penyebut $D(s)$ dapat difaktorkan menjadi perkalian faktor linier real yang saling berbeda:
\begin{equation}
    D(s) = (s - p_1)(s - p_2)\cdots(s - p_n)
\end{equation}
maka $F(s)$ dapat didekomposisi menjadi:
\begin{equation}
    F(s) = \frac{A_1}{s - p_1} + \frac{A_2}{s - p_2} + \dots + \frac{A_n}{s - p_n}
    \label{eq:pfe_distinct}
\end{equation}
Koefisien residu $A_k$ dapat dihitung secara cepat menggunakan \textbf{Metode Penutupan Heaviside} (\textit{Heaviside Cover-up Method}):
\begin{equation}
    A_k = \lim_{s \to p_k} \left[ (s - p_k) F(s) \right] = \left. \frac{N(s)}{\frac{d}{ds}D(s)} \right|_{s = p_k}
    \label{eq:heaviside_coverup}
\end{equation}
Dengan mengambil invers Laplace suku demi suku:
\begin{equation}
    f(t) = \mathcal{L}^{-1}\{F(s)\} = \sum_{k=1}^{n} A_k e^{p_k t} \, u(t)
    \label{eq:inv_distinct}
\end{equation}

\subsection{Kasus 2: Kutub-Kutub Real Berulang (Repeated Real Poles)}
Jika penyebut memiliki kutub berulang sebanyak $m$ kali pada $s = p_1$, yaitu faktor $(s - p_1)^m$, maka suku dekomposisi untuk kutub tersebut memiliki $m$ pecahan berderajat naik:
\begin{equation}
    F(s) = \frac{B_m}{(s - p_1)^m} + \frac{B_{m-1}}{(s - p_1)^{m-1}} + \dots + \frac{B_1}{s - p_1} + \sum \text{suku kutub lainnya}
    \label{eq:pfe_repeated}
\end{equation}
Koefisien residu tertinggi $B_m$ diperoleh langsung dengan metode penutupan biasa:
\begin{equation}
    B_m = \lim_{s \to p_1} \left[ (s - p_1)^m F(s) \right]
\end{equation}
Sedangkan koefisien derajat yang lebih rendah dihitung dengan menggunakan operasi diferensiasi berantai:
\begin{equation}
    B_{m-k} = \frac{1}{k!} \lim_{s \to p_1} \frac{d^k}{ds^k} \left[ (s - p_1)^m F(s) \right], \quad k = 1, 2, \dots, m-1
    \label{eq:repeated_diff_formula}
\end{equation}
Mengingat pasangan invers transformasi baku:
\begin{equation}
    \mathcal{L}^{-1}\left\{ \frac{1}{(s - p_1)^{k+1}} \right\} = \frac{t^k}{k!} e^{p_1 t} u(t)
\end{equation}
maka kemunculan kutub ganda selalu merepresentasikan sinyal waktu yang berinteraksi dengan polinomial waktu (misalnya $t e^{-\alpha t}$ pada sistem redaman kritis / \textit{critically damped}).

\subsection{Kasus 3: Kutub Kompleks Konjugat (Complex Conjugate Poles)}
Sistem rangkaian osilator, filter, dan jaringan redaman kurang (\textit{underdamped}) memiliki pasangan kutub kompleks konjugat $s = -\alpha \pm j\omega_d$. Faktor kuadratik tak tereduksi pada penyebut berbentuk:
\begin{equation}
    D_k(s) = (s + \alpha)^2 + \omega_d^2
\end{equation}
Bentuk pecahan parsial yang paling stabil secara numerik dan intuitif adalah membentuk pembilang berformat kombinasi sinus dan kosinus teredam:
\begin{equation}
    F_{\text{sub}}(s) = \frac{A(s + \alpha) + B \omega_d}{(s + \alpha)^2 + \omega_d^2}
    \label{eq:complex_pfe_standard}
\end{equation}
Menggunakan Teorema Pergeseran Pertama (translasi frekuensi $s \to s + \alpha$):
\begin{equation}
    \mathcal{L}^{-1}\left\{ \frac{s + \alpha}{(s + \alpha)^2 + \omega_d^2} \right\} = e^{-\alpha t} \cos(\omega_d t) u(t)
\end{equation}
\begin{equation}
    \mathcal{L}^{-1}\left\{ \frac{\omega_d}{(s + \alpha)^2 + \omega_d^2} \right\} = e^{-\alpha t} \sin(\omega_d t) u(t)
\end{equation}
Sehingga invers waktu fisisnya adalah:
\begin{equation}
    f_{\text{sub}}(t) = e^{-\alpha t} \left[ A \cos(\omega_d t) + B \sin(\omega_d t) \right] u(t)
    \label{eq:complex_inv_time}
\end{equation}

\section{Pemodelan Fisis Rangkaian Listrik di Domain \texorpdfstring{$s$}{s}}

Keunggulan terbesar Transformasi Laplace dalam rekayasa elektro adalah kemampuannya mentransformasikan elemen pasif rangkaian ($R, L, C$) beserta memori energi awalnya menjadi \textbf{komponen aljabar setara di domain $s$}.

\subsection{Resistor Linier}
Pada resistor ideal, hubungan tegangan dan arus murni seketika tanpa efek memori:
\begin{equation}
    v_R(t) = R \, i_R(t) \quad \xrightarrow{\quad \mathcal{L} \quad} \quad V_R(s) = R \, I_R(s)
\end{equation}
Impedansi operasional resistor adalah $Z_R(s) = R$, dan admitansinya adalah $Y_R(s) = 1/R = G$.

\subsection{Induktor dengan Arus Awal}
Hubungan konstitutif induktor dalam domain waktu adalah $v_L(t) = L \frac{di_L}{dt}$. Mengaplikasikan transformasi Laplace pada operasi turunan:
\begin{equation}
    V_L(s) = \mathcal{L}\left\{ L \frac{di_L}{dt} \right\} = L \left[ s I_L(s) - i_L(0^-) \right]
\end{equation}
Persamaan ini dapat dituliskan dalam dua model ekivalen yang sangat penting:
\begin{enumerate}[leftmargin=*]
    \item \textbf{Model Seri (Ekivalen Thevenin):}
    \begin{equation}
        V_L(s) = sL \cdot I_L(s) - L \, i_L(0^-)
        \label{eq:inductor_thevenin}
    \end{equation}
    Model ini memandang induktor sebagai sebuah impedansi pasif $Z_L(s) = sL$ yang terhubung seri dengan sumber tegangan semu impulsif bernilai konstan $L \, i_L(0^-)$ dengan polaritas yang melawan kenaikan arus.
    \item \textbf{Model Paralel (Ekivalen Norton):}
    Menyelesaikan persamaan untuk $I_L(s)$:
    \begin{equation}
        I_L(s) = \frac{V_L(s)}{sL} + \frac{i_L(0^-)}{s}
        \label{eq:inductor_norton}
    \end{equation}
    Model ini memandang induktor sebagai admitansi pasif $Y_L(s) = \frac{1}{sL}$ yang terhubung paralel dengan sumber arus semu bernilai $\frac{i_L(0^-)}{s}$ yang mengalir searah dengan arus mula-mula.
\end{enumerate}

\subsection{Kapasitor dengan Tegangan Awal}
Hubungan konstitutif kapasitor dalam domain waktu adalah $i_C(t) = C \frac{dv_C}{dt}$. Menerapkan transformasi Laplace:
\begin{equation}
    I_C(s) = C \left[ s V_C(s) - v_C(0^-) \right]
\end{equation}
Serupa dengan induktor, kita mendapatkan dua representasi sirkuit domain $s$:
\begin{enumerate}[leftmargin=*]
    \item \textbf{Model Paralel (Ekivalen Norton):}
    \begin{equation}
        I_C(s) = sC \cdot V_C(s) - C \, v_C(0^-)
        \label{eq:capacitor_norton}
    \end{equation}
    Kapasitor dimodelkan sebagai admitansi pasif $Y_C(s) = sC$ yang terhubung paralel dengan sumber arus semu konstan impulsif $C \, v_C(0^-)$ yang meninggalkan terminal positif.
    \item \textbf{Model Seri (Ekivalen Thevenin):}
    Menyelesaikan persamaan untuk $V_C(s)$:
    \begin{equation}
        V_C(s) = \frac{1}{sC} \cdot I_C(s) + \frac{v_C(0^-)}{s}
        \label{eq:capacitor_thevenin}
    \end{equation}
    Kapasitor dimodelkan sebagai impedansi pasif $Z_C(s) = \frac{1}{sC}$ yang terhubung seri dengan sumber tegangan semu berundak $\frac{v_C(0^-)}{s}$ yang polaritasnya searah dengan orientasi tegangan mula-mula.
\end{enumerate}

Representasi visual sirkuit ekivalen domain $s$ disajikan pada Gambar~\ref{fig:s_domain_models}.

\begin{figure}[htbp]
\centering
\begin{tikzpicture}[scale=0.95, every node/.style={transform shape}]
    % Induktor Thevenin
    \draw (0,2.5) to[L, l={$sL$}] (2.5,2.5)
          to[battery2, l_={$L i(0^-)$}] (4.5,2.5) -- (5.0,2.5);
    \node[above] at (2.5,3.3) {\small\textbf{Induktor: Model Seri (Thevenin)}};
    
    % Kapasitor Thevenin
    \draw (6.5,2.5) to[C, l={$\frac{1}{sC}$}] (9.0,2.5)
          to[battery2, l_={$\frac{v(0^-)}{s}$}] (11.0,2.5) -- (11.5,2.5);
    \node[above] at (9.0,3.3) {\small\textbf{Kapasitor: Model Seri (Thevenin)}};
    
    % Induktor Norton
    \draw (0,0) -- (1.5,0) to[L, l={$sL$}] (1.5,-2) -- (0,-2);
    \draw (1.5,0) -- (3.5,0) to[I, l_={$\frac{i(0^-)}{s}$}] (3.5,-2) -- (1.5,-2);
    \draw (3.5,0) -- (5.0,0);
    \draw (3.5,-2) -- (5.0,-2);
    \node[above] at (2.5,0.4) {\small\textbf{Induktor: Model Paralel (Norton)}};
    
    % Kapasitor Norton
    \draw (6.5,0) -- (8.0,0) to[C, l={$sC$}] (8.0,-2) -- (6.5,-2);
    \draw (8.0,0) -- (10.0,0) to[I, l_={$C v(0^-)$}] (10.0,-2) -- (8.0,-2);
    \draw (10.0,0) -- (11.5,0);
    \draw (10.0,-2) -- (11.5,-2);
    \node[above] at (9.0,0.4) {\small\textbf{Kapasitor: Model Paralel (Norton)}};
\end{tikzpicture}
\caption{Model ekivalen Thevenin dan Norton elemen penyimpan energi di domain frekuensi kompleks $s$.}
\label{fig:s_domain_models}
\end{figure}

Tabel~\ref{tab:s_domain_summary} merangkum seluruh parameter elemen rangkaian di domain waktu dan domain frekuensi kompleks $s$.

\begin{table}[htbp]
\centering
\caption{Ringkasan impedansi, admitansi, dan sumber semu elemen rangkaian di domain $s$.}
\label{tab:s_domain_summary}
\small
\begin{tabularx}{\linewidth}{l >{\centering\arraybackslash}X >{\centering\arraybackslash}X >{\centering\arraybackslash}X >{\centering\arraybackslash}X}
\toprule
\textbf{Elemen} & \textbf{Hubungan $t$} & \textbf{Impedansi $Z(s)$} & \textbf{Sumber Seri Thevenin} & \textbf{Sumber Paralel Norton} \\
\midrule
Resistor ($R$) & $v = R i$ & $R$ & $0$ & $0$ \\
Induktor ($L$) & $v = L \frac{di}{dt}$ & $sL$ & $-L i(0^-)$ & $+\frac{i(0^-)}{s}$ \\
Kapasitor ($C$) & $i = C \frac{dv}{dt}$ & $\frac{1}{sC}$ & $+\frac{v(0^-)}{s}$ & $-C v(0^-)$ \\
\bottomrule
\end{tabularx}
\end{table}

\section{Solusi Rangkaian RLC Dinamis \& Dekomposisi Respon ZIR dan ZSR}

Ketika kita menyusun persamaan simpul (KCL) atau persamaan jerat (KVL) langsung di domain $s$ menggunakan model ekivalen Gambar~\ref{fig:s_domain_models}, variabel luaran sistem, sebutlah $Y(s)$, secara alami terurai menjadi dua komponen ortogonal:
\begin{equation}
    Y(s) = \underbrace{Y_{\text{ZIR}}(s)}_{\text{Respon Masukan-Nol}} + \underbrace{Y_{\text{ZSR}}(s)}_{\text{Respon Keadaan-Nol}}
    \label{eq:zir_zsr_split}
\end{equation}
di mana:
\begin{itemize}[leftmargin=*]
    \item \textbf{Respon Masukan-Nol (\textit{Zero-Input Response} / ZIR):} Respon transien murni yang bersumber semata-mata dari energi awal yang terperangkap di dalam induktor dan kapasitor ($i(0^-) \ne 0$, $v(0^-) \ne 0$), dengan mengasumsikan seluruh sumber eksitasi luar dinonaktifkan (\textit{deactivated} / $X(s) = 0$).
    \item \textbf{Respon Keadaan-Nol (\textit{Zero-State Response} / ZSR):} Respon sistem yang dibangkitkan secara eksklusif oleh sumber masukan luar $X(s)$, dengan mengasumsikan seluruh elemen penyimpan energi berada dalam kondisi istirahat mula-mula tanpa energi awal ($i(0^-) = 0$, $v(0^-) = 0$).
\end{itemize}

Pada kondisi ZSR, hubungan antara masukan $X(s)$ dan luaran $Y_{\text{ZSR}}(s)$ dihubungkan oleh sebuah fungsi rasional fundamental yang disebut \textbf{Fungsi Transfer} (\textit{Transfer Function}) $H(s)$:
\begin{equation}
    H(s) \triangleq \left. \frac{Y(s)}{X(s)} \right|_{\text{kondisi awal} = 0}
    \label{eq:transfer_function}
\end{equation}
Invers Laplace dari fungsi transfer merupakan respon impuls dari sistem:
\begin{equation}
    h(t) = \mathcal{L}^{-1}\{H(s)\}
\end{equation}
Sehingga respon sistem terhadap masukan sebarang $x(t)$ di domain waktu setara dengan integral konvolusi:
\begin{equation}
    y_{\text{ZSR}}(t) = x(t) * h(t) = \int_{0}^{t} x(\tau) h(t - \tau) \, d\tau
\end{equation}

Transformasi Laplace berhasil mengubah konvolusi integral domain waktu yang sangat rumit menjadi sekadar perkalian aljabar skalar di domain frekuensi kompleks: $Y_{\text{ZSR}}(s) = H(s) X(s)$.

\section{Teorema Nilai Awal (IVT) dan Teorema Nilai Akhir (FVT)}

Sering kali dalam analisis praktis, rekayasawan tidak memerlukan profil riwayat waktu transien secara detail untuk seluruh rentang waktu $0 \le t < \infty$. Kita sering kali hanya ingin mengetahui:
\begin{enumerate}[leftmargin=*]
    \item Berapakah lonjakan arus seketika tepat setelah sakelar ditutup ($t = 0^+$)?
    \item Berapakah nilai tegangan akhir pada keadaan tunak setelah seluruh proses transien meluruh sempurna ($t \to \infty$)?
\end{enumerate}
Untuk menjawab pertanyaan tersebut tanpa harus melakukan ekspansi pecahan parsial dan invers Laplace yang panjang, kita menggunakan \textbf{Teorema Nilai Awal} (\textit{Initial Value Theorem}) dan \textbf{Teorema Nilai Akhir} (\textit{Final Value Theorem}).

\subsection{Teorema Nilai Awal (Initial Value Theorem / IVT)}

\begin{theorem}[Teorema Nilai Awal]
Jika fungsi $f(t)$ dan turunan pertamanya $f'(t)$ memiliki transformasi Laplace $F(s)$, serta limit $\lim_{s \to \infty} s F(s)$ ada, maka nilai fungsi sesaat tepat setelah $t = 0$ ($t = 0^+$) diberikan oleh:
\begin{equation}
    f(0^+) = \lim_{t \to 0^+} f(t) = \lim_{s \to \infty} s F(s)
    \label{eq:ivt_theorem}
\end{equation}
\end{theorem}

\begin{proof}
Mulai dari transformasi Laplace dari turunan pertama:
\begin{equation*}
    \mathcal{L}\{f'(t)\} = \int_{0^-}^{\infty} f'(t) e^{-st} \, dt = s F(s) - f(0^-)
\end{equation*}
Ambil limit $s \to \infty$ pada kedua ruas:
\begin{equation*}
    \lim_{s \to \infty} \left[ s F(s) - f(0^-) \right] = \lim_{s \to \infty} \int_{0^-}^{\infty} f'(t) e^{-st} \, dt
\end{equation*}
Karena $e^{-st} \to 0$ saat $s \to \infty$ untuk seluruh $t > 0$, integral di ruas kanan mendekati nilai integral pada selang infinitesimal $t \in [0^-, 0^+]$:
\begin{equation*}
    \lim_{s \to \infty} \int_{0^-}^{\infty} f'(t) e^{-st} \, dt = \int_{0^-}^{0^+} f'(t) \, dt = f(0^+) - f(0^-)
\end{equation*}
Substitusi hasil ini ke persamaan semula:
\begin{equation*}
    \lim_{s \to \infty} s F(s) - f(0^-) = f(0^+) - f(0^-) \implies \lim_{s \to \infty} s F(s) = f(0^+) \qedhere
\end{equation*}
\end{proof}

\subsection{Teorema Nilai Akhir (Final Value Theorem / FVT)}

\begin{theorem}[Teorema Nilai Akhir]
Jika fungsi $f(t)$ dan turunannya memiliki transformasi Laplace $F(s)$, dan \textbf{seluruh kutub dari fungsi $s F(s)$ terletak secara ketat di setengah bidang kiri} (\textit{Left-Half Plane} / LHP, $\text{Re}\{s\} < 0$), maka nilai keadaan tunak (\textit{steady-state}) saat $t \to \infty$ diberikan oleh:
\begin{equation}
    f(\infty) = \lim_{t \to \infty} f(t) = \lim_{s \to 0} s F(s)
    \label{eq:fvt_theorem}
\end{equation}
\end{theorem}

\begin{proof}
Mulai kembali dari definisi transformasi turunan:
\begin{equation*}
    \lim_{s \to 0} \left[ s F(s) - f(0^-) \right] = \lim_{s \to 0} \int_{0^-}^{\infty} f'(t) e^{-st} \, dt
\end{equation*}
Dengan menukar limit dan integral (dijamin valid karena syarat stabilitas kutub LHP):
\begin{equation*}
    \int_{0^-}^{\infty} f'(t) \left( \lim_{s \to 0} e^{-st} \right) \, dt = \int_{0^-}^{\infty} f'(t) \, dt = \lim_{t \to \infty} f(t) - f(0^-) = f(\infty) - f(0^-)
\end{equation*}
Sehingga:
\begin{equation*}
    \lim_{s \to 0} s F(s) - f(0^-) = f(\infty) - f(0^-) \implies \lim_{s \to 0} s F(s) = f(\infty) \qedhere
\end{equation*}
\end{proof}

\begin{tcolorbox}[colback=red!5!white,colframe=red!75!black,title=\textbf{Peringatan Kritis Aplikasi Teorema Nilai Akhir (FVT)}]
\textbf{Perhatian:} Syarat bahwa seluruh kutub dari $sF(s)$ berada di $\text{Re}\{s\} < 0$ adalah \textbf{syarat mutlak}. Jika fungsi $f(t)$ mengandung osilasi sinusoidal murni tak teredam (kutub di sumbu imajiner $s = \pm j\omega$) atau sinyal tak stabil yang meledak (kutub di $\text{Re}\{s\} > 0$), maka limit $\lim_{t \to \infty} f(t)$ secara fisis \textbf{tidak ada}. Menerapkan rumus $\lim_{s \to 0} s F(s)$ secara membabi buta pada kasus tersebut akan menghasilkan nilai numerik palsu yang menyesatkan!
\end{tcolorbox}

\section{Studi Kasus Industri: Transient Recovery Voltage (TRV) pada Pemutus Tenaga Gardu Induk 150 kV PLN}

\subsection{Konteks Masalah Rekayasa Transien Tenaga Listrik}
Pada sistem transmisi tenaga listrik tegangan tinggi 150 kV PLN, salah satu manuver paling kritis adalah pemutusan arus gangguan hubung singkat oleh Pemutus Tenaga (\textit{Circuit Breaker} / CB). Ketika kontak mekanis pemutus tenaga mulai terpisah, terbentuk busur api listrik (\textit{electric arc}) yang sangat panas di antara kedua kontak. Pemadaman busur api hanya dapat terjadi saat arus bolak-balik melintasi titik nol alami (\textit{natural current zero}).

Tepat pada detik busur api padam, tegangan sistem di sisi sumber transmisi tidak serta-merta kembali ke nilai nominal $50\text{ Hz}$ yang tenang. Karena adanya induktansi ekivalen transformator dan saluran ($L_s$) serta kapasitansi parasitik busbar dan bushing ($C_s$), terjadi fenomena peralihan frekuensi tinggi yang sangat dahsyat yang disebut \textbf{Transient Recovery Voltage (TRV)}. Fenomena ini diatur ketat dalam standar internasional \textbf{IEC 62271-100} dan \textbf{IEEE Std C37.011}.

\subsection{Formulasi Rangkaian Transien TRV di Domain \texorpdfstring{$s$}{s}}
Rangkaian ekivalen satu fasa saat pemutusan arus hubung singkat dapat dimodelkan sebagai rangkaian $LC$ paralel teredam oleh resistansi kebocoran $R_p$, yang dieksitasi oleh tegangan sumber bolak-balik puncak $V_m$ saat arus nol ($t = 0$):
\begin{equation}
    v_{\text{sumber}}(t) = V_m \cos(\omega_0 t) \approx V_m, \quad \text{untuk selang waktu mikrodetik transien}
\end{equation}
Beban induktif murni menyebabkan arus tertinggal $90^\circ$ di belakang tegangan, sehingga saat arus melintasi nol, tegangan sumber berada tepat pada puncaknya ($V_m$).

Tegangan di seberang celah kontak pemutus tenaga $V_{\text{TRV}}(s)$ memenuhi persamaan pembagi tegangan operasional:
\begin{equation}
    V_{\text{TRV}}(s) = V_m \cdot \frac{Z_p(s)}{s L_s + Z_p(s)} \cdot \frac{1}{s}
\end{equation}
di mana $Z_p(s) = \frac{1}{s C_s}$. Tanpa redaman resistif ($R_p \to \infty$):
\begin{equation}
    V_{\text{TRV}}(s) = \frac{V_m}{s} \cdot \frac{\frac{1}{s C_s}}{s L_s + \frac{1}{s C_s}} = \frac{V_m}{s} \cdot \frac{1}{s^2 L_s C_s + 1} = \frac{V_m \omega_n^2}{s (s^2 + \omega_n^2)}
    \label{eq:trv_laplace}
\end{equation}
di mana $\omega_n = \frac{1}{\sqrt{L_s C_s}}$ adalah frekuensi alami osilasi transien (biasanya berkisar antara puluhan kHz hingga ratusan kHz).

\subsection{Inversi Analitis \& Laju Kenaikan Tegangan (RRRV)}
Melakukan ekspansi pecahan parsial pada Persamaan~\eqref{eq:trv_laplace}:
\begin{equation}
    \frac{\omega_n^2}{s (s^2 + \omega_n^2)} = \frac{A}{s} + \frac{B s + C}{s^2 + \omega_n^2}
\end{equation}
Dengan metode penutupan Heaviside:
\begin{equation*}
    A = \left. \frac{\omega_n^2}{s^2 + \omega_n^2} \right|_{s = 0} = 1
\end{equation*}
Menyamakan penyebut:
\begin{equation*}
    \omega_n^2 = 1(s^2 + \omega_n^2) + (Bs + C)s = (1 + B)s^2 + Cs + \omega_n^2
\end{equation*}
Diperoleh $B = -1$ dan $C = 0$. Sehingga bentuk domain $s$ adalah:
\begin{equation}
    V_{\text{TRV}}(s) = V_m \left[ \frac{1}{s} - \frac{s}{s^2 + \omega_n^2} \right]
\end{equation}
Mengaplikasikan invers Transformasi Laplace:
\begin{equation}
    v_{\text{TRV}}(t) = V_m \left[ 1 - \cos(\omega_n t) \right] u(t)
    \label{eq:trv_solution_time}
\end{equation}

\textbf{Interpretasi Rekayasa Listrik:}
Persamaan~\eqref{eq:trv_solution_time} mengungkap dua bahaya mematikan bagi pemutus tenaga gardu induk:
\begin{enumerate}[leftmargin=*]
    \item \textbf{Tegangan Puncak Dua Kali Lipat (\textit{Double Peak Amplitude}):} Pada saat $\omega_n t = \pi$ ($t = \pi \sqrt{L_s C_s}$), nilai kosinus bernilai $-1$, sehingga tegangan transien melonjak hingga:
    \begin{equation}
        v_{\text{TRV, maks}} = V_m [1 - (-1)] = 2 V_m
    \end{equation}
    Insulasi dielektrik media gas $\text{SF}_6$ atau ruang hampa (\textit{vacuum}) pada celah kontak harus mampu menahan tegangan dua kali lipat tegangan sistem puncak ini tanpa mengalami tembus tembus listrik (\textit{dielectric breakdown}).
    \item \textbf{Laju Kenaikan Tegangan Pemulihan (\textit{Rate of Rise of Recovery Voltage} / RRRV):}
    \begin{equation}
        \text{RRRV} = \left. \frac{dv_{\text{TRV}}}{dt} \right|_{\text{maks}} = V_m \omega_n = \frac{V_m}{\sqrt{L_s C_s}}
    \end{equation}
    Jika RRRV melampaui laju pemulihan kekuatan dielektrik celah kontak pemutus tenaga ($dV_{\text{diel}}/dt$), maka busur api akan menyala kembali (\textit{re-ignition} atau \textit{re-strike}), yang dapat memicu ledakan katastrofik pada kubikel pemutus tenaga gardu induk PLN!
\end{enumerate}

\section{Contoh Soal Terhitung Numerik Langkah demi Langkah}

\begin{example}[Inversi Pecahan Parsial dengan Kutub Real Berulang \& Berbeda]
Tentukan invers transformasi Laplace $y(t) = \mathcal{L}^{-1}\{Y(s)\}$ dari fungsi aljabar rasional berikut:
\begin{equation*}
    Y(s) = \frac{2s^2 + 10s + 11}{(s + 1)(s + 2)^2}
\end{equation*}
\end{example}

\noindent\textbf{Penyelesaian Langkah demi Langkah:}
\begin{enumerate}[leftmargin=*]
    \item \textbf{Identifikasi Derajat dan Pemfaktoran Kutub:}
    Derajat pembilang adalah 2 dan penyebut adalah 3 ($\text{deg}(N) < \text{deg}(D)$), sehingga fungsi merupakan pecahan sejati. Terdapat satu kutub real sederhana pada $s = -1$ dan satu kutub real ganda ($m=2$) pada $s = -2$.
    
    \item \textbf{Formulasi Bentuk Pecahan Parsial:}
    \begin{equation*}
        Y(s) = \frac{A}{s + 1} + \frac{B_2}{(s + 2)^2} + \frac{B_1}{s + 2}
    \end{equation*}
    
    \item \textbf{Perhitungan Koefisien Residu:}
    \begin{itemize}[leftmargin=*]
        \item Untuk kutub sederhana $s = -1$ (metode penutupan Heaviside):
        \begin{align*}
            A &= \left. (s + 1) Y(s) \right|_{s = -1} = \left. \frac{2s^2 + 10s + 11}{(s + 2)^2} \right|_{s = -1} \\
              &= \frac{2(-1)^2 + 10(-1) + 11}{(-1 + 2)^2} = \frac{3}{1} = \mathbf{3}
        \end{align*}
        \item Untuk kutub ganda derajat tertinggi $B_2$ pada $s = -2$:
        \begin{align*}
            B_2 &= \left. (s + 2)^2 Y(s) \right|_{s = -2} = \left. \frac{2s^2 + 10s + 11}{s + 1} \right|_{s = -2} \\
                &= \frac{2(-2)^2 + 10(-2) + 11}{-2 + 1} = \frac{-1}{-1} = \mathbf{1}
        \end{align*}
        \item Untuk koefisien derajat pertama $B_1$ pada $s = -2$ (menggunakan turunan):
        \begin{align*}
            B_1 &= \left. \frac{d}{ds} \left[ (s + 2)^2 Y(s) \right] \right|_{s = -2} = \left. \frac{d}{ds} \left( \frac{2s^2 + 10s + 11}{s + 1} \right) \right|_{s = -2} \\
            &= \left. \frac{(4s + 10)(s + 1) - (2s^2 + 10s + 11)(1)}{(s + 1)^2} \right|_{s = -2} \\
            &= \frac{[4(-2) + 10](-2 + 1) - [2(-2)^2 + 10(-2) + 11]}{(-2 + 1)^2} \\
            &= \frac{(2)(-1) - (-1)}{(-1)^2} = \frac{-2 + 1}{1} = \mathbf{-1}
        \end{align*}
    \end{itemize}
    
    \item \textbf{Rekonstruksi Invers Domain Waktu:}
    Menuliskan kembali $Y(s)$:
    \begin{equation*}
        Y(s) = \frac{3}{s + 1} + \frac{1}{(s + 2)^2} - \frac{1}{s + 2}
    \end{equation*}
    Menerapkan invers transformasi Laplace suku demi suku:
    \begin{equation*}
        y(t) = \mathbf{\left( 3 e^{-t} + t e^{-2t} - e^{-2t} \right) u(t)} \quad \qed
    \end{equation*}
\end{enumerate}

\begin{example}[Solusi Lengkap Transien RLC Seri dengan Energi Awal di Domain $s$]
Sebuah rangkaian RLC seri memiliki parameter: resistansi $R = 4\,\Omega$, induktansi $L = 1\text{ H}$, dan kapasitansi $C = 0{,}2\text{ F}$. Rangkaian memiliki kondisi awal arus induktor $i(0^-) = 2\text{ A}$ dan tegangan kapasitor $v_C(0^-) = 8\text{ V}$. Pada $t = 0$, rangkaian dihubungkan secara tiba-tiba ke sumber tegangan searah konstan $V_{\text{dc}} = 20\text{ V}$. Tentukan persamaan analitis arus transien $i(t)$ untuk $t \ge 0$ menggunakan analisis domain $s$.
\end{example}

\noindent\textbf{Penyelesaian Langkah demi Langkah:}
\begin{enumerate}[leftmargin=*]
    \item \textbf{Konversi Parameter ke Nilai Impedansi Domain $s$:}
    \begin{itemize}[leftmargin=*]
        \item Resistor: $Z_R(s) = R = 4\,\Omega$
        \item Induktor: $Z_L(s) = sL = s(1) = s$, dengan sumber tegangan semu Thevenin berlawanan arah: $L i(0^-) = 1 \times 2 = 2\text{ V}$
        \item Kapasitor: $Z_C(s) = \frac{1}{sC} = \frac{1}{0{,}2s} = \frac{5}{s}$, dengan sumber tegangan semu Thevenin searah: $\frac{v_C(0^-)}{s} = \frac{8}{s}$
        \item Sumber masukan luar tangga: $V_{\text{in}}(s) = \mathcal{L}\{20 u(t)\} = \frac{20}{s}$
    \end{itemize}
    
    \item \textbf{Penyusunan Persamaan Hukum Kirchoff Tegangan (KVL) di Domain $s$:}
    Menelusuri satu jerat tertutup searah jarum jam:
    \begin{equation*}
        V_{\text{in}}(s) - Z_R(s) I(s) - \left[ Z_L(s) I(s) - L i(0^-) \right] - \left[ Z_C(s) I(s) + \frac{v_C(0^-)}{s} \right] = 0
    \end{equation*}
    Substitusi nilai numerik:
    \begin{equation*}
        \frac{20}{s} - 4 I(s) - [s I(s) - 2] - \left[ \frac{5}{s} I(s) + \frac{8}{s} \right] = 0
    \end{equation*}
    
    \item \textbf{Manipulasi Aljabar untuk Memisahkan Arus $I(s)$:}
    Kelompokkan seluruh suku yang mengandung $I(s)$ di ruas kiri:
    \begin{equation*}
        I(s) \left( s + 4 + \frac{5}{s} \right) = \frac{20}{s} - \frac{8}{s} + 2 = \frac{12}{s} + 2 = \frac{2s + 12}{s}
    \end{equation*}
    Kalikan kedua ruas dengan $s$:
    \begin{equation*}
        (s^2 + 4s + 5) I(s) = 2s + 12 \implies I(s) = \frac{2s + 12}{s^2 + 4s + 5}
    \end{equation*}
    
    \item \textbf{Analisis Kutub dan Penyelesaian Kompleks Konjugat:}
    Akar-akar dari penyebut $s^2 + 4s + 5 = 0$:
    \begin{equation*}
        s_{1,2} = \frac{-4 \pm \sqrt{16 - 20}}{2} = -2 \pm j 1
    \end{equation*}
    Sistem berada pada kondisi \textbf{redaman kurang} (\textit{underdamped}) dengan faktor atenuasi $\alpha = 2$ dan frekuensi sudut teredam $\omega_d = 1\text{ rad/s}$.
    
    Melengkapkan kuadrat pada penyebut:
    \begin{equation*}
        s^2 + 4s + 5 = (s + 2)^2 + 1^2
    \end{equation*}
    Tulis pembilang dalam kombinasi basis $(s + 2)$ dan $\omega_d = 1$:
    \begin{equation*}
        2s + 12 = 2(s + 2) + 8 = 2(s + 2) + 8(1)
    \end{equation*}
    Maka ekspresi $I(s)$ menjadi:
    \begin{equation*}
        I(s) = 2 \cdot \frac{s + 2}{(s + 2)^2 + 1^2} + 8 \cdot \frac{1}{(s + 2)^2 + 1^2}
    \end{equation*}
    
    \item \textbf{Invers Laplace Domain Waktu:}
    \begin{equation*}
        i(t) = \mathbf{e^{-2t} \left[ 2 \cos(t) + 8 \sin(t) \right] u(t)\text{ [A]}} \quad \qed
    \end{equation*}
    
    \item \textbf{Verifikasi Cepat Menggunakan Teorema Nilai Awal (IVT):}
    \begin{equation*}
        i(0^+) = \lim_{s \to \infty} s I(s) = \lim_{s \to \infty} \frac{s(2s + 12)}{s^2 + 4s + 5} = \lim_{s \to \infty} \frac{2s^2 + 12s}{s^2 + 4s + 5} = 2\text{ A}
    \end{equation*}
    Hasil ini identik dengan syarat kontinuitas arus induktor $i(0^+) = i(0^-) = 2\text{ A}$.
\end{enumerate}

\begin{example}[Verifikasi Respon Keadaan Tunak dengan FVT \& Invers Sinyal Berpulsa]
Suatu sistem kendali posisi motor servo di industri memiliki fungsi transfer loop tertutup:
\begin{equation*}
    H(s) = \frac{\Theta(s)}{V_{\text{in}}(s)} = \frac{10}{s^2 + 3s + 10}
\end{equation*}
Sistem dieksitasi oleh pulsa tegangan persegi dengan amplitudo $5\text{ V}$ dan durasi $T = 2\text{ detik}$, yaitu $v_{\text{in}}(t) = 5[u(t) - u(t - 2)]$.
\begin{enumerate}[leftmargin=*]
    \item Tentukan nilai posisi akhir keadaan tunak $\theta(\infty)$ jika masukan berupa tangga konstan $v_{\text{in}}(t) = 5 u(t)$ menggunakan FVT.
    \item Tentukan solusi analitis posisi $\theta(t)$ terhadap pulsa durasi terbatas tersebut.
\end{enumerate}
\end{example}

\noindent\textbf{Penyelesaian Langkah demi Langkah:}
\begin{enumerate}[leftmargin=*]
    \item \textbf{Evaluasi Keadaan Tunak Menggunakan FVT:}
    Untuk masukan tangga $V_{\text{in}}(s) = \frac{5}{s}$:
    \begin{equation*}
        \Theta(s) = H(s) V_{\text{in}}(s) = \frac{10}{s^2 + 3s + 10} \cdot \frac{5}{s} = \frac{50}{s(s^2 + 3s + 10)}
    \end{equation*}
    Kutub-kutub dari penyebut adalah $s = 0$ dan akar dari $s^2 + 3s + 10 = 0$, yaitu $s = -1{,}5 \pm j 2{,}7839$. Karena seluruh kutub non-nol memiliki bagian riil negatif ($\text{Re}\{s\} = -1{,}5 < 0$), Teorema Nilai Akhir dapat diaplikasikan:
    \begin{equation*}
        \theta(\infty) = \lim_{s \to 0} s \Theta(s) = \lim_{s \to 0} \left[ s \cdot \frac{50}{s(s^2 + 3s + 10)} \right] = \frac{50}{0 + 0 + 10} = \mathbf{5\text{ radian}}
    \end{equation*}
    
    \item \textbf{Solusi Respon Waktu terhadap Pulsa Terbatas:}
    Didefinisikan respon terhadap tangga satuan $g(t) = \mathcal{L}^{-1}\left\{ \frac{H(s)}{s} \right\}$.
    Faktorkan suku kuadratik: $(s + 1{,}5)^2 + (\sqrt{7{,}75})^2$, dengan $\omega_d = \sqrt{7{,}75} \approx 2{,}784\text{ rad/s}$.
    \begin{equation*}
        \frac{10}{s[(s + 1{,}5)^2 + 7{,}75]} = \frac{1}{s} - \frac{s + 3}{(s + 1{,}5)^2 + 7{,}75} = \frac{1}{s} - \frac{(s + 1{,}5) + 1{,}5}{(s + 1{,}5)^2 + 7{,}75}
    \end{equation*}
    Invers tangga satuan menghasilkan:
    \begin{equation*}
        g(t) = \left[ 1 - e^{-1{,}5 t} \left( \cos(2{,}784 t) + \frac{1{,}5}{2{,}784} \sin(2{,}784 t) \right) \right] u(t)
    \end{equation*}
    Karena masukan pulsa adalah $v_{\text{in}}(t) = 5 u(t) - 5 u(t - 2)$, berdasarkan sifat linieritas dan Teorema Pergeseran Waktu Kedua:
    \begin{equation*}
        \theta(t) = \mathbf{5 g(t) - 5 g(t - 2) u(t - 2)} \quad \qed
    \end{equation*}
\end{enumerate}

\section{Praktikum Komputasi Numerik Terbuka Berbasis Julia}

Pada bagian ini, kita mengimplementasikan simulasi komputasi numerik ilmiah berbasis bahasa pemrograman modern \textbf{Julia} (versi 1.12+). Program ini memodelkan dinamika transien rangkaian RLC berenergi awal (Contoh Soal 9.2) menggunakan pustaka \texttt{DifferentialEquations.jl}, membandingkan hasil simulasi numerik adaptif terhadap kurva analitis hasil Invers Transformasi Laplace, serta memverifikasi konvergensi Teorema Nilai Awal dan Akhir.

\begin{lstlisting}[style=juliastyle, caption={Skrip Julia untuk simulasi transien RLC dan validasi Invers Laplace.}, label={lst:julia_rlc_laplace}]
# -------------------------------------------------------------
# MODUL 7: SIMULASI TRANSIEN RLC DAN VALIDASI INVERS LAPLACE
# Persamaan Diferensial 2026 - Teknik Elektro UNIB
# Pengampu: Ir. Novalio Daratha, Ph.D. & M. Arfan, M.T.
# -------------------------------------------------------------

using DifferentialEquations
using Plots

# 1. Definisi Parameter Fisik Rangkaian RLC
const R = 4.0      # Resistansi (Ohm)
const L = 1.0      # Induktansi (Henry)
const C = 0.2      # Kapasitansi (Farad)
const Vdc = 20.0   # Tegangan Sumber Tangga (Volt)

# Kondisi Awal Energi Tersimpan
const i0 = 2.0     # Arus Awal Induktor iL(0-) (Ampere)
const vc0 = 8.0    # Tegangan Awal Kapasitor vC(0-) (Volt)

# 2. Pemodelan Ruang Keadaan (State-Space) Domain Waktu
# State vector u = [i_L, v_C]
# du[1]/dt = di_L/dt = (Vdc - R*i_L - v_C) / L
# du[2]/dt = dv_C/dt = i_L / C
function rlc_transient!(du, u, p, t)
    iL = u[1]
    vC = u[2]
    du[1] = (Vdc - R * iL - vC) / L
    du[2] = iL / C
end

# Rentang waktu simulasi
tspan = (0.0, 5.0)
u0 = [i0, vc0]

# Pembentukan dan Penyelesaian Masalah Nilai Awal (IVP)
prob = ODEProblem(rlc_transient!, u0, tspan)
sol = solve(prob, Tsit5(), reltol=1e-8, abstol=1e-10)

# 3. Solusi Analitis Eksak Berbasis Invers Transformasi Laplace
# i(t) = exp(-2t) * [2*cos(t) + 8*sin(t)]
i_exact(t) = exp(-2.0 * t) * (2.0 * cos(t) + 8.0 * sin(t))

# 4. Evaluasi Numerik dan Verifikasi IVT/FVT
t_eval = range(0.0, 5.0, length=500)
i_num = [sol(t)[1] for t in t_eval]
i_ana = [i_exact(t) for t in t_eval]

# Verifikasi Nilai Awal dan Akhir
ivt_val = i_exact(0.0)
println("=== HASIL EVALUASI TEOREMA NILAI AWAL & AKHIR ===")
println("Nilai Awal Terhitung i(0+): ", ivt_val, " A (Eksak: 2.0 A)")
println("Nilai Akhir t=5s        i(5):  ", i_exact(5.0), " A (Mendekati 0 A)")

# 5. Visualisasi Grafik Komparasi Presisi Tinggi
plot(t_eval, i_num, label="Simulasi Numerik (Tsit5)", lw=2.5, color=:blue,
     xlabel="Waktu t [detik]", ylabel="Arus Induktor i(t) [Ampere]",
     title="Dinamika Transien Arus RLC: Numerik vs Invers Laplace",
     legend=:topright, grid=true, size=(800, 480))
plot!(t_eval, i_ana, label="Solusi Analitis Invers Laplace", lw=1.5,
      linestyle=:dash, color=:red)

# Simpan hasil grafik
savefig("simulasi_transien_rlc_modul7.png")
println("Simulasi sukses selesai. Grafik disimpan sebagai simulasi_transien_rlc_modul7.png")
\end{lstlisting}

\section{Evaluasi \& Bank Soal Terstruktur Taksonomi Bloom (C1 s.d. C6)}

Untuk menguji penguasaan materi Modul Minggu 7 secara mendalam dan komprehensif, mahasiswa diwajibkan menyelesaikan bank soal evaluasi berikut:

\begin{enumerate}[leftmargin=*]
    \item \textbf{C1 (Mengingat):} 
    Sebutkan tiga jenis kutub penyebut dalam dekomposisi pecahan parsial fungsi aljabar rasional $F(s)$, serta tuliskan bentuk umum komponen fungsi waktu $f(t)$ yang dihasilkan oleh masing-masing jenis kutub tersebut!
    
    \item \textbf{C2 (Memahami):} 
    Jelaskan secara fisis mengapa sumber semu pada model Thevenin induktor di domain $s$ berwujud sumber tegangan konstan impulsif $L i(0^-)$, sedangkan pada kapasitor berwujud sumber tegangan berundak $v_C(0^-)/s$. Hubungkan penjelasan Anda dengan konsep kontinuitas fluks magnetik dan muatan listrik!
    
    \item \textbf{C3 (Menerapkan):} 
    Diberikan fungsi aljabar rasional pada domain frekuensi kompleks $s$:
    \begin{equation*}
        V(s) = \frac{4s^2 + 16s + 15}{(s + 3)(s^2 + 4s + 13)}
    \end{equation*}
    Hitung invers transformasi Laplace $v(t) = \mathcal{L}^{-1}\{V(s)\}$ menggunakan metode ekspansi pecahan parsial secara lengkap dan teliti!
    
    \item \textbf{C4 (Menganalisis):} 
    Sebuah induktor $L = 0{,}5\text{ H}$ dengan arus awal $i(0^-) = 4\text{ A}$ dihubungkan paralel dengan kapasitor $C = 0{,}05\text{ F}$ yang memiliki tegangan awal $v_C(0^-) = 20\text{ V}$ pada $t = 0$. Tidak ada sumber luar yang terhubung.
    \begin{enumerate}
        \item Gambarkan rangkaian ekivalen domain $s$ menggunakan model admitansi paralel (Norton).
        \item Analisis dan tentukan ekspresi tegangan transien bersama $v(t)$ pada kedua terminal elemen untuk $t \ge 0$.
    \end{enumerate}
    
    \item \textbf{C5 (Mengevaluasi):} 
    Suatu sistem dinamik memiliki luaran $Y(s) = \frac{s + 5}{s^2 + 25}$. Seorang mahasiswa menerapkan Teorema Nilai Akhir dan menyimpulkan bahwa $y(\infty) = \lim_{s \to 0} s Y(s) = \lim_{s \to 0} \frac{s(s+5)}{s^2 + 25} = 0$. Evaluasi keabsahan kesimpulan tersebut dan jelaskan kesalahan fundamental yang dilakukan mahasiswa tersebut ditinjau dari posisi kutub sistem!
    
    \item \textbf{C6 (Menciptakan/Merancang):} 
    Pada gardu induk 150 kV PLN dengan induktansi hubung singkat $L_s = 20\text{ mH}$ dan kapasitansi liar busbar $C_s = 50\text{ nF}$, pemutus tenaga mengalami lonjakan TRV yang sangat berbahaya dengan frekuensi osilasi transien tinggi.
    Rancanglah nilai resistansi peredam snubber $R_s$ yang harus dipasang paralel dengan kontak pemutus tenaga agar osilasi transien TRV teredam secara kritis (\textit{critically damped}), sehingga menghilangkan tegangan lebih puncak dua kali lipat ($2 V_m$)!
\end{enumerate}

\section{Rangkuman, Glosarium Istilah Teknis, dan Referensi Standar}

\subsection{Rangkuman Inti Perkuliahan}
\begin{enumerate}[leftmargin=*]
    \item Invers Transformasi Laplace $\mathcal{L}^{-1}\{F(s)\}$ merupakan gerbang rekonstruksi dari solusi aljabar domain frekuensi kompleks $s$ kembali ke ranah waktu fisis $t$.
    \item Metode Ekspansi Pecahan Parsial (PFE) mengubah rasio polinomial rasional $N(s)/D(s)$ menjadi penjumlahan suku-suku elementer yang kutub-kutubnya berkorespondensi langsung dengan spektrum sinyal fisis (eksponensial riil, polinomial teredam, atau gelombang sinus/kosinus teredam).
    \item Pemodelan rangkaian di domain $s$ mengkonversi hukum kalkulus diferensial-integral menjadi hukum aljabar linier murni. Memori energi awal kapasitor dan induktor terakomodasi secara otomatis sebagai sumber semu Thevenin atau Norton.
    \item Respon sistem terdekomposisi secara alami menjadi Respon Masukan-Nol (ZIR) akibat energi awal internal dan Respon Keadaan-Nol (ZSR) akibat eksitasi sumber luar.
    \item Teorema Nilai Awal (IVT) dan Nilai Akhir (FVT) memberikan pintasan analitis yang sangat ampuh untuk menginspeksi nilai ekstrem $t \to 0^+$ dan $t \to \infty$ secara langsung tanpa perlu melakukan inversi pecahan parsial secara penuh.
\end{enumerate}

\subsection{Glosarium Istilah Teknis Rekayasa}
\begin{description}[leftmargin=!,labelwidth=3.5cm]
    \item[Bromwich Contour] Garis kontur integrasi kompleks pada bidang $s$ yang sejajar sumbu imajiner dan terletak di sebelah kanan seluruh singularitas.
    \item[Cover-up Method] Teknik penutupan faktor penyebut oleh Oliver Heaviside untuk menghitung koefisien residu pecahan parsial secara langsung.
    \item[ZIR / ZSR] \textit{Zero-Input Response} (respon tanpa sumber luar) dan \textit{Zero-State Response} (respon tanpa energi awal).
    \item[Fictitious Source] Sumber tegangan atau arus semu yang merepresentasikan syarat awal memori energi pada elemen induktor atau kapasitor.
    \item[TRV] \textit{Transient Recovery Voltage}; tegangan peralihan berfrekuensi tinggi yang muncul di seberang celah kontak pemutus tenaga sesaat setelah arus busur api padam.
    \item[RRRV] \textit{Rate of Rise of Recovery Voltage}; laju kenaikan tegangan pemulihan transien per satuan waktu, dinyatakan dalam kV/$\mu$s.
\end{description}

\subsection{Referensi Standar Industri \& Akademis}
\begin{enumerate}[leftmargin=*]
    \item Zill, D. G. (2018). \textit{A First Course in Differential Equations with Modeling Applications} (11th ed.). Cengage Learning.
    \item Kreyszig, E. (2011). \textit{Advanced Engineering Mathematics} (10th ed.). John Wiley \& Sons.
    \item Alexander, C. K., \& Sadiku, M. N. O. (2021). \textit{Fundamentals of Electric Circuits} (7th ed.). McGraw-Hill Education.
    \item Hayt, W. H., Kemmerly, J. E., \& Durbin, S. M. (2019). \textit{Engineering Circuit Analysis} (9th ed.). McGraw-Hill Education.
    \item IEEE Std C37.011-2019. \textit{IEEE Guide for the Application of Transient Recovery Voltage for AC High-Voltage Circuit Breakers}. IEEE Power and Energy Society.
    \item IEC 62271-100:2021. \textit{High-voltage switchgear and controlgear -- Part 100: Alternating-current circuit-breakers}. International Electrotechnical Commission.
\end{enumerate}

\end{document}
